\documentclass{article}
\usepackage[a4paper,scale=0.8,hcentering,bindingoffset=8mm]{geometry} % A4纸大小，缩放80%，设置奇数页右边留空多一点
\usepackage{hyperref}      % 超链接
\usepackage{listings}      % 代码块
\usepackage{courier}       % 字体
\usepackage{fontspec}      % 字体
\usepackage{fancyhdr}      % 页眉页脚相关宏包
\usepackage{lastpage}      % 引用最后一页
\usepackage{amsmath,amsthm,amsfonts,amssymb,bm} %数学
\usepackage{graphicx}      % 图片
\usepackage{subcaption}    % 图片描述
\usepackage{longtable,booktabs} % 表格
\usepackage{ctex}
\usepackage{soul}
\usepackage{dot2texi}
\usepackage{tikz}
\usepackage{multirow}
% \setmainfont{Noto Sans Kaithi}
\usetikzlibrary{automata, positioning, arrows}
\lstset{                  %设置代码块
         basicstyle=\footnotesize\ttfamily,% 基本风格
         numbers=left,    % 行号
         numbersep=10pt,  % 行号间隔 
         tabsize=4,       % 缩进
         extendedchars=true, % 扩展符号？
         breaklines=true, % 自动换行
         language=C,
         frame=leftline,  % 框架左边竖线
         xleftmargin=5pt,% 竖线左边间距
         showspaces=false,% 空格字符加下划线
         showstringspaces=false,% 字符串中的空格加下划线
         showtabs=false,  % 字符串中的tab加下划线
 }
\pagestyle{fancy}         % 页眉页脚风格
\fancyhf{}                % 清空当前设置
\fancyfoot[C]{\thepage\ / \pageref{LastPage}}%页脚中间显示 当前页 / 总页数，把\label{LastPage}放在最后
\begin{document} 
    \begin{titlepage}       % 封面
        \centering
        \includegraphics[width=\textwidth]{../SUSTC_LOGO.png}
        % \vspace*{\baselineskip}
        \rule{\textwidth}{1.6pt}\vspace*{-\baselineskip}\vspace*{2pt}
        \rule{\textwidth}{0.4pt}\\[\baselineskip]
        {\LARGE COMPILIER @Liu Yepang 2019\\[\baselineskip]\small for SUSTech CSE}
        \\[0.2\baselineskip]
        \rule{\textwidth}{0.4pt}\vspace*{-\baselineskip}\vspace{3.2pt}
        \rule{\textwidth}{1.6pt}\\[\baselineskip]
        \scshape
        \vspace*{\baselineskip}
        {\Large HomeWork 3\par }
        Edited by \\[\baselineskip] {汪至圆\par}
        {\Large 11610634\par }
        \vfill
        {\scshape 2019} \\{\large SHENZHEN}\par
    \end{titlepage}
    \section{Exercise 1 (Top-Down Parsing): }
    Consider the following grammar G:
    $$
        \begin{array}{c}{S \rightarrow a B} \\ {B \rightarrow S+B | \epsilon}\end{array}
    $$
        \subsection{Construct the predictive parsing table for G. [20 points]}
            $$
                \begin{array}{c}
                    {FIRST(S):\{a\}}\\
                    {FIRST(B):\{a,\epsilon\}}\\
                    {FELLOW(S)\{+,\$\}}\\
                    {FELLOW(B)\{+,\$\}}
                \end{array}
            $$
            \begin{table}[!htbp]
                \centering
                \caption{Parsing Table}
                \label{tab:aStrangeTable}
                \begin{tabular}{|c|c|c|c|}
                    \hline
                    \multirow{2}*{Non-Terminal} & \multicolumn{3}{c|}{INPUT SYMBOL}\\
                    \cline{2-4}
                    ~ & a & + & \$\\
                    \hline
                    S & $S \rightarrow a B$ & & \\
                    B & $B \rightarrow S+B$ & $B \rightarrow \epsilon$ & $B \rightarrow \epsilon$ \\
                    \hline
                \end{tabular}
            \end{table}
        \subsection{Is the grammar LL(1)? [10 points]}
            For the $B \rightarrow S+B | \epsilon$
            \begin{itemize}
                \item ${FIRST(S+B)=\{a\}}$ and ${FIRST(\epsilon)=\{\epsilon\}}$, $FIRST(S+B)\bigcup FIRST(\epsilon)=\emptyset$
                \item ${FIRST(\epsilon)=\{\epsilon\}}$, and $FIRST(S+B)\bigcup FELLOW(B)=\emptyset$
            \end{itemize}
            So $B \rightarrow S+B | \epsilon$ is LL(1)
        \subsection{Can an LL(1) parser accept the input string aaaa+++? If yes, please list the moves made by the parser; otherwise, state the reason. Before parsing, please resolve conflicts in the parsing table if any. [20 points]}
\end{document}